% ==============================================================================
%  CHAPTER 3 — METHODOLOGY
% ==============================================================================
\section{Methodology}
\label{sec:method}

I followed a build-and-evaluate approach. I assembled a dataset of first-person video
(Section~\ref{sec:dataset}), built a five-stage pipeline that detects, tracks, classifies
and selectively blurs faces (Sections~\ref{sec:arch} to~\ref{sec:anon}), and evaluated it
against human judgements and against a blanket-blurring baseline
(Section~\ref{sec:evaldesign}).

\subsection{Dataset}
\label{sec:dataset}

\paragraph{Source and selection.}
No public dataset labels which people in first-person video are interacting with the
camera wearer and which are bystanders, so I assembled my own. I collected 21 videos from
publicly available online sources, mainly YouTube; I did not record any footage myself.
A video was included if it was filmed from a first-person viewpoint (a camera worn on the
head or held at eye level), contained faces large enough to be detected, and clearly
belonged to one of two categories:
\begin{itemize}
  \item \textbf{Interaction} (14 videos): the wearer spends time with other people, mostly
        playing card or board games around a table, and in some videos shopping in small
        stores and talking to staff. These videos contain both partners and bystanders,
        for example other customers in a shop.
  \item \textbf{Bystander} (7 videos): the wearer walks through busy outdoor places such
        as street markets and city streets, by day and at night, without talking to
        anyone. Every face in these videos is a bystander.
\end{itemize}
Interaction videos were trimmed to five-minute segments. Bystander videos were used at
their original length of 6 to 78 seconds. Source resolutions ranged from
$1280 \times 720$ to $3840 \times 2160$, and I resized every frame to $1280 \times 720$
before processing so that all thresholds apply to the same frame size. % TODO(Akash): if
% you remember the search terms or channels used to find the videos, add one sentence.

\paragraph{Size of the dataset.}
Counting only videos would understate how much data the system processes, because each
video contains many people and each person appears in many frames. Table~\ref{tab:dataset}
gives the dataset at three levels: videos, frames, and face tracks. A face track is one
person followed through consecutive frames by the tracker, and it is the unit the
classifier makes decisions about. Across the 21 videos the pipeline processed
131{,}187 frames, made 185{,}865 face detections and followed 2{,}736 face tracks, of
which 600 were visible for at least two seconds.

\begin{table}[htbp]
\begin{singlespace}
\centering
\caption{Size of the evaluation dataset at the level of videos, frames and faces.}
\label{tab:dataset}
\renewcommand{\arraystretch}{1.3}
\begin{tabular}{@{}lrrr@{}}
\toprule
 & \textbf{Interaction} & \textbf{Bystander} & \textbf{Total} \\
\midrule
Videos                              & 14         & 7        & 21 \\
Total duration (minutes)            & 70.0       & 2.8      & 72.8 \\
Frames                              & 125{,}986  & 5{,}201  & 131{,}187 \\
Frames containing at least one face & 90{,}398   & 3{,}959  & 94{,}357 \\
Face detections                     & 169{,}247  & 16{,}618 & 185{,}865 \\
Face tracks                         & 2{,}288    & 448      & 2{,}736 \\
Face tracks visible for 2~s or more & 508        & 92       & 600 \\
\bottomrule
\end{tabular}
\end{singlespace}
\end{table}

\paragraph{Balance.}
By duration the dataset is dominated by interaction video (70 of 73 minutes). The
imbalance between partners and bystanders is smaller than this suggests, because the
interaction videos also contain many bystanders: other customers, people at other tables
and passers-by. I handled the remaining imbalance in three ways. First, the labelled
evaluation samples were drawn from every video, so that short bystander videos are
represented alongside long interaction videos. Second, I report bystander protection and
partner preservation as two separate rates (Section~\ref{sec:metrics}), so a large class
cannot hide poor performance on a small one, which would happen with overall accuracy.
Third, I report results separately for the two video categories.

\paragraph{Ground truth.}
Whether a person is interacting with the wearer is a human judgement, so the ground truth
comes from human annotators. I built three labelled sets over the course of the project,
each improving on the previous one. They are described in Section~\ref{sec:groundtruth}
and summarised in Table~\ref{tab:gtsets}. Their accuracy is checked by measuring agreement
between independent annotators (Section~\ref{sec:stats}).

\paragraph{Handling of personal data.}
The videos show members of the public. I kept the videos and all frames derived from them
on my own computer, shared clips only with the annotators for labelling, and did not
publish any frame that shows a face.

\subsection{System Architecture}
\label{sec:arch}

The pipeline processes one frame at a time in five stages (Figure~\ref{fig:pipeline}).
Each frame is passed through detection, tracking, cue extraction, classification and
anonymization, and the output is the same frame with only bystander faces blurred.

\begin{figure}[htbp]
  \centering
  \fbox{\begin{minipage}{0.9\linewidth}\centering\vspace{4pt}
    \textit{Frame} $\rightarrow$ \textbf{Detection} $\rightarrow$ \textbf{Tracking}
    $\rightarrow$ \textbf{Cue Extraction} $\rightarrow$ \textbf{Classification}
    $\rightarrow$ \textbf{Anonymization} $\rightarrow$ \textit{Output}\vspace{4pt}
  \end{minipage}}
  \caption{The five stages of the pipeline.}
  \label{fig:pipeline}
\end{figure}

Identity and appearance are kept separate. The tracker gives each face a number that only
says ``this is the same box as in the previous frame''; the cue stage looks at the face in
the current frame. No stage computes a face embedding or compares faces with each other,
so the system never learns who anyone is, and the track numbers are discarded when the
video ends.

I implemented the pipeline in Python 3.13 with OpenCV 4.13 for video input and image
processing, Ultralytics 8.4 for face detection \cite{yolov8}, my own implementation of
SORT for tracking, and the MediaPipe 0.10 Face Landmarker for facial landmarks
\cite{lugaresi2019mediapipe}. Everything runs on the CPU; no GPU or neural accelerator is
used.

\subsection{Face Detection}
\label{sec:detection}

YOLOv8n-face (Section~\ref{sec:lit_detect}) produces a bounding box and a confidence score for every face.
Each frame is resized to the detector's input of $640 \times 640$ pixels, and overlapping
boxes are merged by non-maximum suppression (NMS). I set the confidence threshold to
$\tau_{\mathrm{det}} = 0.45$ and the NMS overlap threshold to $\tau_{\mathrm{NMS}} = 0.45$
empirically, aiming to suppress low-confidence false detections without losing faces that
are partly hidden or poorly lit. At most 20 faces are kept per frame.

\subsection{Multi-Face Tracking}
\label{sec:tracking}

The cues described below are accumulated over time, so each face needs a stable identity
across frames. I use SORT \cite{sort2016}, which predicts each face box with a
constant-velocity Kalman filter and matches predictions to new detections with the
Hungarian algorithm, using the overlap (IoU) between boxes as the matching cost. A track
that loses its detection is kept for up to 90 frames (3~s at 30~fps), so that a face
briefly hidden behind a hand or another person keeps its identity. A new track is only
passed to the classifier after three consecutive detections, which filters out one-frame
false detections. SORT uses box motion only, so no appearance model of the face is ever
built.

\subsection{Interaction Cues}
\label{sec:cues}

For each tracked face the system computes four cues. I developed them in two phases: the
two geometry cues first (Phase~2a), and the two landmark cues afterwards (Phase~2b).

\subsubsection*{Geometry cues (Phase 2a)}

These cues use only the bounding box. They rely on the centre bias of first-person video:
the wearer tends to face the person they are talking to, so that person appears large and
near the centre of the frame \cite{li2013gaze,fathi2012social}.

\textbf{Proximity.} For a face box $[x_1, y_1, x_2, y_2]$ in a $W \times H$ frame, the
area ratio is
\begin{equation}
  r_{\mathrm{area}} = \frac{(x_2 - x_1)(y_2 - y_1)}{W \cdot H}.
  \label{eq:area_ratio}
\end{equation}
In my videos, a person in conversation with the wearer typically produced values between
0.02 and 0.08, while distant bystanders were usually below 0.003. I set the threshold to
$\tau_{\mathrm{area}} = 0.012$ by looking at the distribution of each track's largest area
ratio across the videos.

\textbf{Centrality.} The distance of the face centre $(\bar{x}, \bar{y})$ from the frame
centre, normalised by the half-diagonal, is
\begin{equation}
  d_{\mathrm{center}} = \frac{\sqrt{(\bar{x} - W/2)^2 + (\bar{y} - H/2)^2}}
                             {\sqrt{(W/2)^2 + (H/2)^2}},
  \label{eq:center_dist}
\end{equation}
which is 0 at the centre and about 1 at the corners. A face counts as central when
$d_{\mathrm{center}} \le \tau_{\mathrm{center}} = 0.45$.

\subsubsection*{Landmark cues (Phase 2b)}

For each tracked face, the box is enlarged by 15\% of its longer side so that the chin and
forehead are included, and the crop is passed to the MediaPipe Face Landmarker
\cite{kartynnik2019facemesh}, which returns 478 facial landmarks. Each crop is processed
on its own, since SORT already handles tracking. If no face is found in the crop, the
landmark cues are left empty for that frame and the classifier falls back to the geometry
cues.

\textbf{Head yaw.} I estimate head pose by solving the Perspective-$n$-Point problem
\cite{lepetit2009epnp} for six landmarks with known positions on a generic 3D face model:
the nose tip, the chin, the outer corners of both eyes and both corners of the mouth
(landmarks 4, 152, 33, 263, 61 and 291). With $\mathbf{P}_i^{3D}$ the model points and
$\mathbf{p}_i^{2D}$ their positions in the crop, the camera matrix is approximated as
\begin{equation}
  K = \begin{bmatrix} f & 0 & c_x \\ 0 & f & c_y \\ 0 & 0 & 1 \end{bmatrix},
  \quad f = W_{\mathrm{crop}}, \quad c_x = W_{\mathrm{crop}}/2, \quad c_y = H_{\mathrm{crop}}/2,
  \label{eq:camera_matrix}
\end{equation}
and OpenCV's iterative solver finds the rotation $\mathbf{r}$ and translation
$\mathbf{t}$ that minimise the reprojection error
\begin{equation}
  \min_{\mathbf{r}, \mathbf{t}} \sum_{i=1}^{6}
  \left\| \mathbf{p}_i^{2D} - \pi\!\left(K,\, \mathbf{r},\, \mathbf{t},\, \mathbf{P}_i^{3D}\right) \right\|^2,
  \label{eq:pnp}
\end{equation}
where $\pi(\cdot)$ is perspective projection. The rotation is converted to a matrix with
the Rodrigues formula and decomposed into Euler angles; the yaw angle
$\theta_{\mathrm{yaw}}$ (turning left or right) is the cue. A face with
$|\theta_{\mathrm{yaw}}| > \tau_{\mathrm{yaw}} = 40^\circ$ is treated as looking away from
the wearer.

\textbf{Mouth aspect ratio (MAR).} Following the eye aspect ratio of Soukupov\'a and
\v{C}ech \cite{soukupova2016}, I measure how far the mouth is open relative to its width.
The upper inner lip $\bar{u}$ is the mean of landmarks 82, 13 and 312, the lower inner lip
$\bar{l}$ the mean of 87, 14 and 317, and the mouth corners are landmarks 78 and 308:
\begin{equation}
  \mathrm{MAR} = \frac{\lVert \bar{u} - \bar{l} \rVert}{\lVert p_{78} - p_{308} \rVert}.
  \label{eq:mar}
\end{equation}
MAR is close to zero when the mouth is closed. In my videos, values above about 0.08
corresponded to an open mouth during speech.

\subsection{Classification}
\label{sec:classification}

The classifier assigns each tracked face one of two labels in every frame:
\textsc{Interactor} (keep visible) or \textsc{Bystander} (blur). It follows three rules
of thumb: when in doubt, blur; never let one noisy frame change a label; and let the
landmark cues veto or support the geometry decision.

\subsubsection*{Decision for a single frame}

For each frame the cues are checked in this order:
\begin{enumerate}
  \item \textbf{Yaw veto.} If landmarks are available and
        $|\theta_{\mathrm{yaw}}| > 40^\circ$, the face is a bystander for this frame.
  \item \textbf{Geometry rule.} The face is an interactor if it is both close and central:
        \begin{equation}
          f_{\mathrm{geo}} = \bigl[r_{\mathrm{area}} \geq \tau_{\mathrm{area}}\bigr]
                             \wedge
                             \bigl[d_{\mathrm{center}} \leq \tau_{\mathrm{center}}\bigr].
          \label{eq:geo_rule}
        \end{equation}
  \item \textbf{Speaking boost.} The system keeps the last $N_w = 30$ MAR values of each
        track and computes the fraction in which the mouth was open,
        \begin{equation}
          s_{\mathrm{vad}} = \frac{1}{|\mathcal{W}|} \sum_{m \in \mathcal{W}}
                             \mathbf{1}\!\left[m \geq \tau_{\mathrm{MAR}}\right],
          \label{eq:vvad}
        \end{equation}
        once at least 10 values are available. If $s_{\mathrm{vad}} \geq
        \tau_{\mathrm{vad}} = 0.30$ and the face passes the proximity test, it is an
        interactor even if it is slightly off-centre, up to
        $1.2\,\tau_{\mathrm{center}}$. This covers a partner who moves to the side while
        speaking.
  \item \textbf{Default.} Otherwise the face is a bystander.
\end{enumerate}

\subsubsection*{Confirmation and latching}

A face changes from \textsc{Bystander} to \textsc{Interactor} only after the single-frame
decision says interactor in $\tau_{\mathrm{enter}} = 2$ consecutive frames. A face that
covers at least 4\% of the frame is confirmed at once, because at that distance the person
is almost certainly face to face with the wearer. Once confirmed, a face stays an
interactor for the rest of the video, even if the cues later fail. Without this latch, a
partner who looked down at their cards or turned to someone else would be blurred and
unblurred repeatedly. The latch accepts a small risk to bystander privacy in exchange for stable
partner preservation. Table~\ref{tab:params} lists all thresholds.

\begin{table}[htbp]
\begin{singlespace}
\centering
\caption{Classification thresholds.}
\label{tab:params}
\renewcommand{\arraystretch}{1.3}
\begin{tabular}{@{}llll@{}}
\toprule
\textbf{Parameter} & \textbf{Symbol} & \textbf{Value} & \textbf{Cue} \\
\midrule
Minimum area ratio            & $\tau_{\mathrm{area}}$   & $0.012$     & Proximity \\
Maximum centre distance       & $\tau_{\mathrm{center}}$ & $0.45$      & Centrality \\
Instant-confirmation area     & $r_{\mathrm{instant}}$   & $0.040$     & Proximity \\
Confirmation streak           & $\tau_{\mathrm{enter}}$  & 2 frames    & Hysteresis \\
Maximum yaw                   & $\tau_{\mathrm{yaw}}$    & $40^\circ$  & Head pose \\
Open-mouth threshold          & $\tau_{\mathrm{MAR}}$    & $0.08$      & Speaking \\
Speaking window               & $N_w$                    & 30 frames   & Speaking \\
Speaking fraction             & $\tau_{\mathrm{vad}}$    & $0.30$      & Speaking \\
\bottomrule
\end{tabular}
\end{singlespace}
\end{table}

\subsection{Selective Anonymization and Baseline}
\label{sec:anon}

Every face labelled \textsc{Bystander} is blurred with a $51 \times 51$ Gaussian kernel
over its box enlarged by 10\% on each side, so that hair and ears are covered too. Faces
labelled \textsc{Interactor} are left untouched. The pipeline writes three videos from the
same detections and tracks: the selective output, a \emph{baseline} in which every
detected face is blurred, and a debug video that draws each decision. Because the
baseline uses the same detector and tracker, any difference between the two outputs comes
from the classification stage alone.

\subsection{Offline Second Pass for Recorded Video}
\label{sec:secondpass}

The latch only works forward in time, so a partner is blurred until the system has seen
enough evidence. When the whole recording is available, two further rules can be applied
after the first pass, using only the decisions already made:
\begin{enumerate}
  \item \textbf{Retroactive latch.} A track confirmed as a partner at any point is kept
        visible for its whole time on screen, including the frames before confirmation.
  \item \textbf{Group rule.} A track that is visible for at least 10~seconds, and is on
        screen together with a confirmed partner for at least half of that time, is
        treated as part of the wearer's group and kept visible. Friends sitting at the same
        table stay in view for minutes, while passers-by rarely stay for 10~seconds.
        Requiring a confirmed partner to be present stops the rule from preserving
        strangers who simply stay in one place, such as a shop assistant.
\end{enumerate}
I designed these rules after studying the errors of the video-clip evaluation
(Section~\ref{sec:results_clips}), so their results are reported separately from the main
evaluation.

\subsection{Evaluation Design}
\label{sec:evaldesign}

\subsubsection*{Ground-truth sets}
\label{sec:groundtruth}

Table~\ref{tab:gtsets} summarises the three labelled sets.

\begin{table}[htbp]
\begin{singlespace}
\centering
\caption{Labelled sets used for evaluation. Each data point is one person (face track)
judged as an interaction partner or a bystander.}
\label{tab:gtsets}
\renewcommand{\arraystretch}{1.3}
\fitwidth{\begin{tabular}{@{}lllcc@{}}
\toprule
\textbf{Set} & \textbf{What annotators saw} & \textbf{Annotators} &
\textbf{Data points} & \textbf{Partners / bystanders} \\
\midrule
A: Phase 2a tracks  & Full videos             & 1 (myself)        & 66  & 26 / 40 \\
B: Still frames     & 10 frames, faces numbered & 4 (incl.\ myself) & 29  & 7 / 22 \\
C: Video clips      & 7~s clips with sound    & 5 (incl.\ myself) & 102 & 69 / 33 \\
\bottomrule
\end{tabular}}
\end{singlespace}
\end{table}

\textbf{Set A.} For the first evaluation I watched each video in full and labelled 66 face
tracks across the 21 videos. % TODO(Akash): add how the 66 tracks were chosen.
A track was labelled a partner if the person had a face-to-face conversation with the
wearer at any point, matching the latching behaviour of the classifier. These labels
could not be reused for Phase~2b: SORT numbers tracks in order of appearance, so when the
detections changed slightly between the Phase~2a and Phase~2b runs, the track numbers
shifted and no longer referred to the same people.

\textbf{Set B.} For Phase~2b I first ran a blind study on still frames. One frame was
chosen automatically from each of the first five bystander and first five interaction
videos: the frame with the most clearly visible faces (box at least 50~pixels on each
side, detection confidence at least 0.50, outside the first 30 frames). Each face was
marked with a plain number, with nothing indicating the system's decision. Four annotators
labelled each face as a partner (``actively talking to the camera wearer'') or a
bystander, or left it blank if they could not decide.

\textbf{Set C.} Still frames make it hard to tell who is talking, so I ran a larger study
with video. From every evaluation video I sampled tracks without looking at the system's
decisions: the four longest tracks of each bystander video, and the four longest plus two
random tracks of each interaction video, using tracks visible for at least 60 frames with
a median face size of at least 32~pixels. Each sampled track became a seven-second clip:
one second of a still frame marking the target face, then six seconds of the original,
unblurred video with its sound, with a yellow box on the target face. The 104 clips were
shown in a random order on a web page. For this study an \emph{interaction partner} was
defined as any member of the wearer's group or conversation, including a friend at the
same table who is not speaking at that moment, and a \emph{bystander} as anyone outside
that group. Five annotators labelled every clip. Two clips turned out not to show a person
at all (a face printed on a magazine and a product package) and were excluded, leaving
102.

In both studies the final label of each face is the majority vote, with ties counted as
bystander, the privacy-safe choice.

\subsubsection*{System decision for each data point}

For Sets A and B the system's decision for a track is \emph{partner} if the track was
confirmed as an interactor at any point in the video, which is how the latch behaves in the
output. For Set C the decision is what a viewer of the output clip actually sees: a face
counts as kept visible if it was unblurred in at least half of the frames in which it was
detected during the clip. I checked these decisions against the rendered output videos by
comparing each face region with the original and the blanket-blurred frames, and all 104
matched.

\subsubsection*{Metrics}
\label{sec:metrics}

I treat \emph{interaction partner} as the positive class, so a true positive (TP) is a
partner kept visible, a false negative (FN) a partner blurred, a false positive (FP) a
bystander left visible, and a true negative (TN) a bystander blurred. The two main
metrics are
\begin{equation}
  \mathrm{PSR} = \frac{TN}{TN + FP}, \qquad \mathrm{IPR} = \frac{TP}{TP + FN}.
  \label{eq:psr_ipr}
\end{equation}
The Privacy Success Rate (PSR) is the share of bystanders who are blurred, and the
Interaction Preservation Rate (IPR) is the share of partners who stay visible. The names
are mine, but the measures follow BystandAR \cite{bystandar}, which reports the percentage
of bystanders protected and the percentage of intended subjects kept accessible. Using the
same two measures makes my results directly comparable with that system. In standard
classification terms, PSR is the specificity and IPR is the recall (sensitivity) of the
partner class.

I report the two separately rather than a single accuracy figure for two reasons. The two
errors have different costs: a false positive exposes someone who did not consent, while a
false negative only removes useful content. And accuracy would be dominated by whichever
class is larger in a given sample. I also report precision and F$_1$, for comparison with
general classification work. Every data point is a person rather than a frame, because the
question an annotator answers, and the decision the latch makes, is about a person.

\subsubsection*{Statistical analysis}
\label{sec:stats}

\textbf{Annotator agreement.} The ground truth is only useful if different people judge
the same faces the same way. Raw percentage agreement overstates this, because two
annotators who both label most faces as bystanders will agree often by chance. Cohen's
kappa corrects for chance agreement \cite{cohen1960}:
\begin{equation}
  \kappa = \frac{p_o - p_e}{1 - p_e},
  \label{eq:kappa}
\end{equation}
where $p_o$ is the observed agreement and $p_e$ the agreement expected if both annotators
labelled at random with their own label frequencies. $\kappa = 1$ means perfect agreement
and $\kappa = 0$ means no better than chance. Cohen's kappa compares two annotators, so I
computed it for every pair. To summarise agreement across all annotators at once, I used
Fleiss' kappa \cite{fleiss1971}, which extends the same idea to any fixed number of
annotators. I interpret values with the scale of Landis and Koch \cite{landis1977}: 0.41
to 0.60 is moderate, 0.61 to 0.80 substantial, and above 0.80 almost perfect agreement.

\textbf{Confidence intervals.} PSR and IPR are proportions estimated from samples of a few
dozen people, so each comes with uncertainty. I report 95\% Wilson score intervals
\cite{wilson1927}. I chose the Wilson interval over the more common normal-approximation
interval because the latter fails exactly where my results lie: for small samples and for
proportions near 0\% or 100\%, it can give intervals below 0\% or above 100\%, or of zero
width when every case is correct. The Wilson interval stays inside $[0, 1]$ and remains
reliable for small samples.

\subsubsection*{Speed and power measurement}

I measured processing time on a laptop with an Intel Core i5-13420H, using only the CPU.
Each stage of the pipeline was timed separately with a high-resolution timer over the
first 300 frames of one interaction video, and the throughput in frames per second (FPS)
is the number of frames divided by the total time. I ran the benchmark with the laptop on
mains power and again on battery, since wearable devices always run on battery and laptops
limit processor speed when unplugged.

\subsubsection*{Projection to edge hardware}

I did not have a Raspberry Pi~5, so I estimated its performance from published
measurements. For the detection stage, which dominates the processing time, I used the
Raspberry Pi~5 benchmarks published by Ultralytics for the YOLOv8n model at the same
$640 \times 640$ input size \cite{ultralytics_rpi}. Of the export formats benchmarked
there, I used two: PyTorch (the format my pipeline currently uses) and NCNN, the fastest
format available on that board. For the remaining stages (frame reading, tracking,
landmarks, classification and blurring), which are ordinary CPU work, I applied a factor
of 3.5. This factor comes from published Geekbench~6 scores, in which the Raspberry Pi~5
scores roughly 2.5 to 3.5 times lower than 13th-generation Intel laptop processors
\cite{geekbench_pi5}; I used the upper end to stay conservative.
